\section{Related Work}
%With advances in GPU technology, training deeper neural networks have become feasible. Popular deeper networks such as VGG and ResNets have produced state-of-the-art performances on standard object recognition, segmentation and face recognition benchmarks. However, training such networks efficiently and robustly is an active area of research.
Many approaches have been proposed to regularize the training procedure of very deep networks. Early stopping and statistical techniques like weight decay are commonly used to prevent overfitting. Specialized techniques such as DropConnect ~\cite{icml2013_wan13}, Dropout ~\cite{srivastava2014dropout} have been successfully applied with very deep networks such as ResNets. Faster convergence of such deep architectures was made possible by Batch Normalization (BN) \cite{ioffe2015batch}. One of the added benefit of BN was that the additional regularization provided during training even made dropout regularization unnecessary in some cases.

The work of Szegedy \emph{et al.} ~\cite{szegedy2013intriguing} showed the existence of adversarial perturbations for computer vision tasks by solving a box-constrained optimization approach to generate these perturbations. They also showed that training the network by feeding back these adversarial examples regularizes the training and makes network resistant to adversarial examples. Due to a relatively expensive layerwise training procedure, their analysis was limited to small datasets and shallow networks.~\cite{goodfellow2014explaining} proposed the fast gradient sign method to generate such adversarial examples. To perform adversarial training, they proposed a modified loss function to also account for loss from adversarial examples. They showed significant improvements in the network's response to adversarial examples and obtained a regularization performance beyond dropout. ~\cite{miyato2017virtual} proposes a virtual adversarial training framework and show its regularization benefits for relatively deep models, while taking three times the normal training time. ~\cite{deepfool} proposed an iterative approach to generate much stronger adversarial perturbations and also presented a score function to measure robustness of classifiers against these examples. Furthermore, recent approaches such as deep contrastive smoothing~\cite{gu2014towards}, distillation  ~\cite{papernot2016distillation} and stability training~\cite{zheng2016improving} have focused solely on improving robustness of the deep models to adversarial inputs. In this work, we present an efficient layerwise approach to adversarial training and demonstrate its ability as a strong regularizer for very deep models beyond the specialized methods mentioned above, in addition to improving model robustness to adversarial inputs.


%~\cite{zheng2016improving} improved the robustness of DNNs against small perturbations caused by cropping, compression and rescaling by adding a stability term in the loss function.

Recent theoretical works such as [\cite{fawzi2015analysis}, ~\cite{fawzi2016robustness}] analyze the effect of random, semi-random and adversarial perturbations on classifier robustness. They presented fundamental upper bounds on the robustness of classifier which depends on factors such as curvature of decision boundary and distinguishability between class cluster centers. Wang \emph{et al.} \cite{WangGQ16} introduce the notion of strong robustness for classifiers and point out that the differences between generalization and robustness by characterizing the topology of the learned classification function. We perform empirical studies that show that the proposed approach improves performance by suppressing those dimensions that are unnecessary for generalization. In addition, we observe that pure adversarial training techniques suppress most dimensions resulting in strong robustness against adversarial examples but less improvement in generalization. To the best of our knowledge, ours is the first work to explore this comparison between regularization and adversarial robustness by providing empirical results on very deep networks. 

% suggest that classifiers are strong robust if the learned feature representation does not include unnecessary feature dimensions and proposed that in order for DNN to be strong robust, one needs to improve feature representation and reduce the effect of unnecessary features. 

%~\cite{szegedy2013intriguing} discussed many interesting properties of these adversarial examples an
%~\cite{understandingAT} proposed a fast gradient sign method to generate such perturbations and presented a modified loss function which minimizes the classification and adversarial loss jointly. ~\cite{foolingimages}


%Neural network trained with these adversarial examples have shown to be robust to such example. It is also observed that these adversarial training provides stronger regularization to the training of deep neural networks. 